\chapter*{About GSET Paper I}
\addcontentsline{toc}{chapter}{About GSET Paper I}
\markboth{About GSET Paper I}{About GSET Paper I}

\section*{The examination}
The \textbf{Gujarat State Eligibility Test (GSET)} is conducted by the GSET Agency, \textbf{The Maharaja Sayajirao University of Baroda, Vadodara}, as the UGC-accredited State Eligibility Test for Gujarat (and Goa in some years). Qualifying GSET makes a candidate eligible for the post of \textbf{Assistant Professor} in universities and colleges of the State.

\section*{Pattern of Paper I}
\begin{center}\small
\begin{tabularx}{\linewidth}{@{}lX@{}}\toprule
\textbf{Feature} & \textbf{Details}\\\midrule
Nature & General paper on teaching and research aptitude, common to all subjects\\
Questions & 50 compulsory MCQs (2017 and earlier: 60 given, any 50 to be answered, the first 50 attempted being evaluated)\\
Marks & 2 marks each, 100 marks in total; no negative marking\\
Medium & Bilingual question booklet (English and Gujarati)\\
Duration & 1 hour (Paper I, followed by Paper II on the subject)\\\bottomrule
\end{tabularx}
\end{center}

\section*{Syllabus (UGC-aligned units)}
\begin{enumerate}[leftmargin=2em]
\item \textbf{Teaching Aptitude}: concept, objectives, levels of teaching (memory, understanding, reflective); learner characteristics; factors affecting teaching; methods of teaching; teaching support systems; evaluation systems.
\item \textbf{Research Aptitude}: meaning, types and characteristics of research; positivism and post-positivism; methods (experimental, descriptive, historical, qualitative, quantitative); steps; thesis and article writing; research ethics.
\item \textbf{Comprehension}: a passage followed by questions.
\item \textbf{Communication}: meaning, types and characteristics; effective communication; verbal and non-verbal; inter-cultural and group communication; classroom communication; barriers; mass media and society.
\item \textbf{Mathematical Reasoning and Aptitude}: types of reasoning; number, letter and code series; relationships; fractions, time and distance, ratio, proportion, percentages, profit and loss, interest, discounting, averages.
\item \textbf{Logical Reasoning}: structure of arguments; deductive and inductive reasoning; analogies; Venn diagrams; Indian logic (pramanas, anumana, vyapti, hetvabhasa).
\item \textbf{Data Interpretation}: sources, acquisition and classification of data; quantitative and qualitative data; graphical representation (bar chart, histogram, pie chart, table chart, line chart); data and governance.
\item \textbf{Information and Communication Technology (ICT)}: general abbreviations and terminology; basics of the internet, intranet, e-mail, audio and video conferencing; digital initiatives in higher education; ICT and governance.
\item \textbf{People, Development and Environment}: development and environment (MDGs and SDGs); human--environment interaction; environmental issues; pollutants and their impact; natural and energy resources; natural hazards and disasters; environmental protection acts and international agreements.
\item \textbf{Higher Education System}: institutions of learning in ancient India; evolution of higher learning after independence; oriental, conventional and non-conventional learning; professional, technical and skill-based education; value education; policies, governance and administration (including NEP~2020).
\end{enumerate}

\section*{Papers in this book}
\begin{center}\small
\begin{tabular}{llcc}\toprule
\textbf{Examination} & \textbf{Date} & \textbf{Booklet} & \textbf{Questions}\\\midrule
November 2025 & 16 Nov 2025 & A & 50\\ December 2024 & 1 Dec 2024 & A & 50\\ November 2023 & 26 Nov 2023 & A & 50\\ November 2022 & 6 Nov 2022 & A & 50\\ December 2021 & 23 Jan 2022 & A & 50\\ December 2019 & 29 Dec 2019 & A & 50\\ September 2018 & 30 Sep 2018 & A & 50\\ August 2017 & 27 Aug 2017 & A & 60\\ September 2016 & 25 Sep 2016 & A & 60\\ October 2014 & Oct 2014 & A & 60\\ December 2013 & Dec 2013 & A & 60\\ September 2013 & Sep 2013 & A & 60\\ October 2011 & Oct 2011 & A & 60\\ October 2010 & Oct 2010 & A & 60\\ December 2008 & 7 Dec 2008 & C & 50\\ July 2006 & Jul 2006 & A & 50\\ July 2004 & 25 Jul 2004 & C & 50\\ December 2003 & 7 Dec 2003 & C & 50\\ December 2002 & 1 Dec 2002 & B & 50\\ January 2002 & 20 Jan 2002 & B & 50\\\midrule
\textbf{Total} & & & \textbf{1{,}070}\\\bottomrule
\end{tabular}
\end{center}
